\section{Introduction}

Single-cell functional genomics screens now make it possible to connect genetic perturbations with transcriptional response phenotypes at scale \citep{dixit2016perturbseq,adamson2016parseq,datlinger2017pooled,xie2017multiplexed,norman2019manifolds,srivatsan2020sciplex,datlinger2021ultra,papalexi2021eccite,frangieh2021multimodal,replogle2022genome,peidli2024scperturb}. These Perturb-seq and related CRISPR screens are increasingly used to study gene function, pathway regulation, combinatorial genetic interactions, and candidate target prioritization, and curated perturbation resources are growing rapidly \citep{perturbase2025,perturbdb2025}. At the same time, computational perturbation-response predictors such as scGen, CPA, GEARS, CellOT, CINEMA-OT, and related single-cell models have made it feasible to infer responses for perturbations or contexts that are only partially observed in the training data \citep{lotfollahi2019scgen,lotfollahi2023cpa,roohani2023gears,bunne2023cellot,dong2023cinemaot}. Broader single-cell representation models and virtual-cell proposals further increase interest in prediction systems that can be interrogated before biological interpretation \citep{lopez2018scvi,gayoso2022scvitools,theodoris2023geneformer,cui2024scgpt,bunne2024virtualcell}. For biological interpretation, however, the key question is not only whether a model has good average accuracy, but whether a specific prediction remains reliable when it is used to prioritize a pathway, perturbation, or follow-up experiment.

Random-split accuracy can obscure this reliability problem. In a random split, train and test data may share perturbations, control strata, cell contexts, support structure, and experimental protocols. Recent perturbation-prediction benchmarks and subtask-decomposition work similarly emphasize that models can behave differently across biologically distinct generalization regimes \citep{generalizable2025benchmark,stamp2024}. Dataset shift, domain generalization, calibration, uncertainty, selective prediction, and conformal prediction literatures all emphasize that average in-distribution accuracy does not guarantee reliable deployment under transfer \citep{zhou2022dgsurvey,koh2021wilds,ovadia2019uncertainty,guo2017calibration,hendrycks2016ood,geifman2017selective,vovk2005algorithmic,angelopoulos2021conformal}. A model can therefore appear strong while failing on the settings where computational prediction is most useful: unseen perturbations, unseen combinations, low-support signatures, new datasets, or external public screens. These failures are biologically consequential because a non-transportable prediction can enter downstream pathway enrichment, perturbation retrieval, or gene-function interpretation as if it were an experimentally supported response. Reliability modeling for perturbation prediction should therefore ask where transfer fails and whether high-risk predictions can be filtered before biological interpretation.

We address this problem with the Perturbation Transportability Layer (PTL), a model-agnostic reliability model for already-generated single-cell perturbation-response predictions. PTL does not replace an expression predictor; it estimates the risk that an already-generated response prediction will fail under transfer. PTL uses four linked components: a stress-split benchmark for perturbation, combination, support, dataset, and external transfer; a relative transportability labeling rule anchored to split-specific difficulty; a deployment-feature keep/abstain model; and a failure-mode atlas that links prediction errors to novelty, support, and dataset boundaries. Using public Perturb-seq screens, we show that random-split winners do not define universal transfer winners, that PTL reduces false transportability compared with confidence filtering, and that severe failures concentrate at biologically interpretable transfer boundaries.
